\documentclass{article}
\usepackage{tikz}
\usetikzlibrary{calc,arrows.meta,positioning,shapes.geometric}
\usepackage{graphicx}
\usepackage{float}
\usepackage{mathptmx}
\usepackage{indentfirst}
\usepackage[utf8]{inputenc}
\usepackage[T5]{fontenc}
\usepackage[paperheight=29.7cm,paperwidth=21cm,right=2cm,left=3cm,top=2cm,bottom=2.5cm]{geometry}
\usepackage[fontsize=13pt]{scrextend}
\usepackage{indentfirst}
\usepackage{lipsum}
\renewcommand{\contentsname}{MỤC LỤC}
\renewcommand{\tablename}{Bảng}
\renewcommand{\listoffigures}{DANH MỤC HÌNH VẼ}
\renewcommand{\listoftables}{DANH MỤC BẢNG BIỂU}
\renewcommand{\figurename}{Hình}
\graphicspath{ {figures/} }
\usepackage{array}
\usepackage{tabularx}
\usepackage{makecell}
\usepackage{hyperref}

\newcolumntype{a}{>{\hsize = 1.1\hsize}X}
\newcolumntype{s}{>{\hsize=2\hsize}X}

\begin{document}

\section{Thu thập dữ liệu (Crawling data)}

Mục tiêu của giai đoạn này là xây dựng một kho ngữ liệu (corpus) gồm các báo cáo
thường niên / báo cáo phát triển bền vững của các doanh nghiệp niêm yết, tập
trung vào nhóm ngành \textit{Xây dựng -- Vật liệu xây dựng -- Bất động sản}. Quy
trình thu thập gồm bốn bước: (1) xác định danh sách doanh nghiệp từ cổng dữ liệu
Vietstock, (2) lập danh sách đường dẫn báo cáo, (3) tải báo cáo tự động, và (4)
giải nén các tập nén để chuẩn hoá đầu ra về dạng tài liệu PDF.

\subsection{Nguồn dữ liệu doanh nghiệp}

Danh sách doanh nghiệp được lấy từ trang tra cứu doanh nghiệp A--Z của Vietstock:
\href{https://finance.vietstock.vn/doanh-nghiep-a-z?page=1}{\texttt{finance.vietstock.vn/doanh-nghiep-a-z}}.
Trang này liệt kê toàn bộ doanh nghiệp niêm yết theo thứ tự bảng chữ cái, kèm mã
chứng khoán, tên công ty và ngành nghề. Từ nguồn này, nhóm nghiên cứu lọc ra các
doanh nghiệp thuộc nhóm ngành mục tiêu và tổng hợp đường dẫn tới báo cáo thường
niên của từng doanh nghiệp theo từng năm vào một tệp Excel danh mục
(\texttt{config/company\_annual\_report.xlsx}).

Tệp Excel được tổ chức theo sheet ngành; sheet được sử dụng có tên
\textit{``Xây dựng -- VLXD -- BĐS''}. Mỗi dòng mô tả một báo cáo và gồm năm cột
chuẩn hoá: mã chứng khoán (\texttt{ma\_ck}), tên công ty (\texttt{ten\_cty}),
tên tài liệu (\texttt{ten\_tl}), năm báo cáo (\texttt{nam}) và đường dẫn tải
(\texttt{url}). Các giá trị mã chứng khoán / tên công ty được điền xuôi
(forward-fill) để xử lý các ô bị gộp (merged cell) trong Excel; những dòng không
có URL hợp lệ (không bắt đầu bằng \texttt{http://} hoặc \texttt{https://}) bị loại bỏ.

\subsection{Tải báo cáo tự động}

Việc tải báo cáo được thực hiện bằng kịch bản
\texttt{crawl\_data/download\_reports.py}. Kịch bản đọc tệp Excel danh mục, sinh
ra cấu trúc thư mục đầu ra có dạng:

\begin{verbatim}
data/raw/annual_report/
    Xây dựng - VLXD - BĐS/
        {Mã CK} - {Tên công ty}/
            {Mã CK}_{Năm}.{ext}     # tệp đơn (pdf, ...)
            {Mã CK}_{Năm}/          # thư mục nếu nguồn là tập nén
\end{verbatim}

Các đặc điểm kỹ thuật chính của bộ tải:

\begin{itemize}
  \item \textbf{Tải song song:} sử dụng \texttt{ThreadPoolExecutor} với 5 luồng
        đồng thời để tăng thông lượng.
  \item \textbf{Chuẩn hoá tên tệp / thư mục:} loại bỏ các ký tự không hợp lệ trên
        Windows (\texttt{<>:"/\textbackslash|?*} và ký tự điều khiển), giới hạn độ
        dài tên để tránh lỗi đường dẫn quá dài.
  \item \textbf{Cho phép chạy lại (resume):} tự động bỏ qua những tệp đã tải thành
        công (kích thước $>0$) hoặc đã được giải nén trước đó, nên có thể dừng và
        chạy lại mà không tải trùng.
  \item \textbf{Cơ chế thử lại:} mỗi URL được thử tối đa 3 lần với thời gian chờ
        tăng dần (backoff tuyến tính theo số lần thử).
  \item \textbf{Giả lập trình duyệt:} đặt header \texttt{User-Agent} kiểu Chrome để
        giảm khả năng bị máy chủ từ chối truy cập.
  \item \textbf{Ghi nhật ký lỗi:} các đường dẫn tải thất bại được ghi vào
        \texttt{download\_log.csv} (kèm mã CK, năm, URL, thông báo lỗi) để rà soát
        và tải bổ sung.
  \item \textbf{Khử trùng lặp đích:} khi cùng một (mã CK, năm) có nhiều URL, tệp
        đích được thêm hậu tố phiên bản (\texttt{\_v2}, \texttt{\_v3}, ...) để
        tránh ghi đè.
\end{itemize}

\subsection{Giải nén và chuẩn hoá định dạng}

Một số báo cáo được công bố dưới dạng tập nén (\texttt{.zip}, \texttt{.rar},
\texttt{.7z}). Bộ tải tích hợp logic giải nén từ
\texttt{crawl\_data/extract\_archives.py}: ngay sau khi tải, mỗi tập nén được
giải vào một thư mục cùng tên (bỏ phần mở rộng) rồi xoá tệp nén gốc, nhằm bảo đảm
đầu ra cuối cùng chỉ còn tài liệu, không còn tập nén.

Đối với tệp \texttt{.zip}, kịch bản xử lý riêng lỗi mã hoá tên tệp tiếng Việt:
các tệp ZIP cũ trên Windows thường mã hoá tên theo \texttt{cp437}, nên tên tệp
được giải mã lại sang \texttt{utf-8} (hoặc \texttt{cp1258} dự phòng) để giữ đúng
dấu tiếng Việt. Tệp \texttt{.rar} và \texttt{.7z} được giải qua công cụ ngoài
(\texttt{UnRAR.exe}, \texttt{7z.exe}). Việc giải nén cũng chạy đa luồng và bỏ qua
các thư mục đích đã tồn tại nội dung.

\medskip
\noindent\textbf{Ghi chú.} Ngoài luồng tải theo danh mục Excel ở trên, kho mã
nguồn còn có \texttt{crawl\_data/crawler.py} -- bộ thu thập chuyên biệt cho một
website doanh nghiệp có bảo vệ Cloudflare. Bộ này dùng \texttt{nodriver}
(Chrome không bị phát hiện tự động hoá), vô hiệu hoá trình xem PDF nội bộ của
Chrome để mỗi lần điều hướng tới một liên kết PDF sẽ kích hoạt tải tệp, và dùng
biểu thức chính quy để lọc đúng các liên kết báo cáo thường niên / ESG. Đây là
phương án bổ trợ khi báo cáo không thể tải trực tiếp qua \texttt{requests}.

\section{Xử lý dữ liệu (Processing data)}

Sau khi có kho PDF, giai đoạn xử lý chuyển corpus thô thành tập câu được
\textit{gán nhãn ESG} bằng một mô hình phân loại ngôn ngữ đã được tinh chỉnh
(fine-tuned). Khác với hướng tiếp cận dựa trên từ điển từ khoá, ở đây
\textbf{mô hình ngôn ngữ chính là bộ phát hiện ESG}: mọi câu trích từ báo cáo đều
được đưa qua mô hình, không qua bất kỳ bước lọc từ khoá hay luật nào. Toàn bộ mã
nằm trong gói \texttt{data\_processing/} và notebook
\texttt{notebooks/kaggle\_esg\_classify.ipynb}.

Pipeline được thiết kế theo mô hình \textit{lai} (hybrid): phần trích xuất văn
bản chạy cục bộ trên CPU, còn phần suy luận mô hình chạy trên GPU (Kaggle). Luồng
tổng thể gồm năm bước (Hình~\ref{fig:pipeline}): (i) trích xuất văn bản PDF,
(ii) tách câu tiếng Việt, (iii) chuẩn bị câu thành tệp JSONL, (iv) phân loại ESG
bằng mô hình ViDeBERTa, và (v) trích xuất các câu đã được gán nhãn ESG.

\begin{figure}[H]
  \centering
  \begin{tikzpicture}[
      node distance=7mm,
      box/.style={draw, rounded corners, align=center, inner sep=4pt,
                  text width=11cm, minimum height=8mm},
      gpu/.style={draw, rounded corners, align=center, inner sep=4pt,
                  text width=11cm, minimum height=8mm, very thick},
      arrow/.style={-{Stealth[length=2.5mm]}, thick}
    ]
    \node[box] (pdf)  {Trích xuất văn bản PDF \\ \footnotesize (\texttt{pdf\_extractor.py} -- PyMuPDF, giữ số trang)};
    \node[box, below=of pdf] (sent) {Tách câu tiếng Việt \\ \footnotesize (\texttt{sentence\_splitter.py} -- \texttt{underthesea})};
    \node[box, below=of sent] (prep) {Chuẩn bị câu $\rightarrow$ JSONL \\ \footnotesize (\texttt{prepare\_sentences.py} -- mọi câu, không lọc ESG)};
    \node[gpu, below=of prep] (clf) {Phân loại ESG bằng mô hình ViDeBERTa \\ \footnotesize (\texttt{esg\_classifier.py} / notebook trên Kaggle GPU)};
    \node[box, below=of clf] (out) {Trích xuất câu đã gán nhãn ESG \\ \footnotesize (\texttt{extract\_esg.py})};
    \draw[arrow] (pdf)  -- (sent);
    \draw[arrow] (sent) -- (prep);
    \draw[arrow] (prep) -- (clf);
    \draw[arrow] (clf)  -- (out);
  \end{tikzpicture}
  \caption{Luồng xử lý trích xuất câu ESG dựa trên mô hình ViDeBERTa. Khối in
  đậm chạy trên GPU (Kaggle); các khối còn lại chạy cục bộ trên CPU.}
  \label{fig:pipeline}
\end{figure}

\subsection{Trích xuất văn bản PDF}

Mô-đun \texttt{pdf\_extractor.py} dùng thư viện PyMuPDF (\texttt{fitz}) để trích
xuất văn bản theo từng trang, giữ nguyên số trang (1-indexed) -- thông tin cần
thiết để truy vết câu về đúng vị trí trong báo cáo. PyMuPDF xử lý tốt dấu tiếng
Việt và bảo toàn thứ tự đọc với bố cục báo cáo thường niên điển hình. Văn bản mỗi
trang được làm sạch: chuẩn hoá Unicode về dạng NFC, loại bỏ ký tự gạch nối mềm
(soft hyphen), ghép lại các từ bị ngắt bởi dấu gạch nối cuối dòng, gộp khoảng
trắng ngang và rút gọn các dòng trống liên tiếp.

\subsection{Tách câu tiếng Việt}

Mô-đun \texttt{sentence\_splitter.py} ưu tiên dùng \texttt{underthesea.sent\_tokenize}
-- bộ tách câu nhận biết tiếng Việt (xử lý đúng các viết tắt như ``TP.HCM'',
``Q.1'', số thập phân ``1.500'', ...); nếu thư viện không khả dụng, hệ thống lùi
về một bộ tách bằng biểu thức chính quy. Quy trình tách câu gồm: (1) chia trang
thành các khối đoạn theo dòng trống, (2) ghép các dòng bị ngắt mềm (do bố cục)
thành văn bản liền mạch trong mỗi khối, (3) chạy bộ tách câu, và (4) lọc bỏ
nhiễu: dòng mục lục (nhận diện qua dãy dấu chấm dẫn trang), số trang đơn lẻ, mảnh
chỉ gồm số/ký hiệu bảng biểu, và các câu quá ngắn (dưới 25 ký tự hoặc dưới 4 từ).

\subsection{Chuẩn bị câu cho mô hình}

Mô-đun \texttt{prepare\_sentences.py} là nửa \textit{cục bộ} của pipeline. Nó duyệt
một tệp PDF hoặc cả thư mục PDF (đệ quy), tái sử dụng đúng hai bước trích xuất và
tách câu ở trên, rồi ghi \textbf{toàn bộ} câu ra một tệp JSONL (kèm bản sao CSV để
mở bằng Excel). Điểm mấu chốt: bước này \textbf{không áp dụng bất kỳ bộ lọc ESG
nào} -- mọi câu đều được giữ lại và chuyển cho mô hình, vì việc phát hiện ESG được
giao hoàn toàn cho mô hình ViDeBERTa ở bước sau. Mỗi dòng JSONL là một bản ghi câu
gồm bốn trường: tên PDF nguồn (\texttt{source\_pdf}), số trang (\texttt{page}),
chỉ số câu trong trang (\texttt{sentence\_index}) và nội dung câu (\texttt{text}).

\subsection{Phân loại ESG bằng mô hình ViDeBERTa}

Đây là nửa \textit{GPU} của pipeline, được hiện thực hoá trong notebook
\texttt{notebooks/kaggle\_esg\_classify.ipynb} (và mô-đun tương đương
\texttt{esg\_classifier.py} có thể chạy độc lập trên CPU/GPU). Mô hình sử dụng là
\texttt{nguyen599/ViDeBERTa-v3-ESG-base} -- một mô hình DeBERTa-v2 đã được tinh
chỉnh cho bài toán phân loại ESG trên văn bản tiếng Việt/tiếng Anh.

\subsubsection{Cơ chế phân loại đa nhãn}

Mô hình có bốn nhãn: \textit{Neutral} (trung tính), \textit{Environmental} (Môi
trường), \textit{Social} (Xã hội) và \textit{Governance} (Quản trị). Do cấu hình
\texttt{problem\_type} là \texttt{multi\_label\_classification}, mỗi nhãn nhận một
điểm xác suất \textbf{độc lập} qua hàm \textit{sigmoid} (không dùng softmax). Lưu
ý kỹ thuật: thẻ cấu hình công bố ghi nhãn trung tính là ``Neural'', nên mã nguồn
chuẩn hoá lại thành ``Neutral''.

\subsubsection{Quy tắc gán nhãn}

Từ vector điểm sigmoid của mỗi câu, hệ thống quyết định theo hai quy tắc tách bạch
(Bảng~\ref{tab:thresholds}):

\begin{itemize}
  \item \textbf{Gán trụ cột E/S/G:} một câu được gắn nhãn một trụ cột khi điểm
        sigmoid của trụ cột đó $\ge$ \texttt{THRESHOLD} (mặc định $0.45$). Một câu
        có thể nhận nhiều trụ cột cùng lúc.
  \item \textbf{Xác định tính liên quan ESG:} một câu được coi là \textit{có liên
        quan ESG} khi điểm \textit{Neutral} $<$ \texttt{NEUTRAL\_THRESHOLD}
        (mặc định $0.5$) -- \textbf{không} căn cứ vào việc một trụ cột đơn lẻ có
        vượt ngưỡng hay không.
\end{itemize}

\noindent Lý do của quy tắc thứ hai: trên thực tế mô hình này phân bổ xác suất
giống softmax (tổng các trụ cột và Neutral xấp xỉ $1$) và hiếm khi gắn hai trụ cột
cùng lúc. Do đó một câu rõ ràng mang nội dung ESG có thể bị \textit{chia tín hiệu}
giữa hai trụ cột (ví dụ câu về tuân thủ luật lao động cho $S=0.44$ và $G=0.44$,
không trụ cột nào vượt $0.45$). Việc neo tính liên quan ESG vào điểm Neutral giúp
thu hồi đúng những câu này, đồng thời vẫn loại được câu khuôn mẫu thuần tuý
(thường có Neutral $\ge 0.95$).

\begin{table}[H]
  \centering
  \caption{Các siêu tham số suy luận của mô hình phân loại ESG.}
  \label{tab:thresholds}
  \begin{tabularx}{\linewidth}{l l X}
    \hline
    \textbf{Tham số} & \textbf{Giá trị} & \textbf{Ý nghĩa} \\
    \hline
    \texttt{THRESHOLD}         & $0.45$ & Ngưỡng sigmoid để gán một trụ cột E/S/G. \\
    \texttt{NEUTRAL\_THRESHOLD} & $0.5$  & Câu liên quan ESG khi điểm Neutral nhỏ hơn ngưỡng này. \\
    \texttt{MAX\_LENGTH}       & $256$  & Độ dài token tối đa (giới hạn mô hình là $512$; câu thường ngắn). \\
    \texttt{BATCH\_SIZE}       & $64$   & Kích thước lô khi suy luận theo lô. \\
    \hline
  \end{tabularx}
\end{table}

\subsubsection{Quy trình vận hành lai (local + Kaggle)}

\begin{enumerate}
  \item Chạy \texttt{prepare\_sentences.py} cục bộ để sinh tệp \texttt{*.jsonl}
        chứa toàn bộ câu.
  \item Tải tệp \texttt{*.jsonl} lên Kaggle dưới dạng \textit{Dataset} và đính kèm
        vào notebook.
  \item Chạy notebook trên Kaggle với \textit{Accelerator = GPU T4} (không dùng
        P100 do ảnh PyTorch đã bỏ hỗ trợ kiến trúc này) và bật \textit{Internet}
        để tải mô hình lần đầu.
  \item Suy luận theo lô: mỗi câu được token hoá (đệm/cắt ở $256$ token), đưa qua
        mô hình, áp sigmoid trên logits để thu điểm bốn nhãn, rồi áp hai quy tắc
        gán nhãn ở trên.
  \item Tải kết quả về: mỗi tệp đầu vào sinh một tệp
        \texttt{<tên>\_classified.jsonl}. Mỗi dòng kết quả hợp nhất bốn trường câu
        gốc với ba trường dự đoán: \texttt{scores} (điểm bốn nhãn), \texttt{labels}
        (danh sách trụ cột vượt ngưỡng) và \texttt{esg} (cờ liên quan ESG).
\end{enumerate}

\subsection{Trích xuất câu đã gán nhãn ESG}

Bước cuối, mô-đun \texttt{extract\_esg.py} đọc các tệp JSONL đã phân loại, giữ lại
những câu có ít nhất một nhãn trụ cột (\texttt{labels} khác rỗng) và tỉa mỗi bản
ghi về các trường cần thiết cho bước xây dựng đồ thị tri thức (GraphRAG) về sau:
\texttt{source\_file}, \texttt{source\_pdf}, \texttt{page}, \texttt{sentence\_index},
\texttt{text}, \texttt{labels} và \texttt{scores}. Mô-đun phản chiếu cấu trúc thư
mục đầu vào và xuất bốn dạng kết quả: tệp \texttt{<tên>\_extracted.jsonl} cho từng
nguồn, tệp gộp \texttt{esg\_all\_records.jsonl}, tệp gom theo tài liệu nguồn
\texttt{esg\_by\_document.json}, và tệp thống kê tổng hợp \texttt{esg\_stats.json}.
Kết quả này chính là tập câu ESG sạch, có truy vết tới đúng tài liệu và trang,
dùng làm đầu vào cho các bước phân tích / xây dựng đồ thị tiếp theo.

\section{Trích xuất chỉ số ESG (KPI extraction)}

Sau khi đã có tập câu được gán nhãn ESG, giai đoạn này chuyển dòng câu \textit{phi
cấu trúc} thành các \textit{quan sát chỉ số} (KPI observations) có cấu trúc, đầy đủ
giá trị, đơn vị, năm và phân loại (mục tiêu / đã đạt / dự báo / kỳ gốc). Đây là
cầu nối giữa phần xử lý ngôn ngữ ở Mục~2 và bước xây dựng đồ thị tri thức tiếp
theo, vì lớp đồ thị được tổ chức quanh nút \texttt{KPIObservation} -- mỗi nút
mang một con số đo lường gắn với tài liệu nguồn. Toàn bộ logic được hiện thực
trong mô-đun \texttt{src/extract\_kpi\_from\_jsonl.py}; bộ từ vựng KPI mà bước
này tham chiếu được xây dựng riêng theo quy trình mô tả trong tài liệu
\texttt{docs/KPI\_DEFINITIONS\_CONSTRUCTION\_BUILD.md}.

\subsection{Lý do cần một bước trích xuất riêng}

Bộ phân loại ViDeBERTa ở Mục~2 chỉ trả lời câu hỏi \textit{``câu này có liên quan
ESG hay không, thuộc trụ cột E/S/G nào''}, nhưng không cho biết \textit{``đang
đo cái gì, bằng bao nhiêu, đơn vị nào, năm nào''}. Trong khi đó, lược đồ đồ thị
yêu cầu mỗi \texttt{KPIObservation} phải có đầy đủ các trường định lượng
(\texttt{kpi\_type}, \texttt{value}, \texttt{unit}, \texttt{kind},
\texttt{direction}, \texttt{year}, \texttt{target\_year}, \texttt{baseline\_year},
\texttt{source\_id}). Vì vậy cần một bước biến đổi: từ câu ESG dạng tự nhiên
sang bản ghi KPI có kiểu xác định, được neo về một mã \texttt{kpi\_type} trong
bộ từ vựng kiểm soát.

Đặc điểm của báo cáo ESG tiếng Việt khiến cách tiếp cận luật/biểu thức chính quy
không phù hợp: số liệu được viết theo nhiều dạng (\texttt{1.500}, \texttt{1,5
triệu}, \texttt{15\%}); đơn vị thường ngầm định; cùng một chỉ số có thể được
diễn đạt rất khác nhau giữa các doanh nghiệp; và một câu có thể chứa đồng thời
\textit{kỳ gốc}, \textit{mục tiêu} và \textit{kết quả đã đạt}. Bước này do đó
sử dụng \textbf{mô hình ngôn ngữ lớn (LLM)} -- cụ thể là \texttt{gemini-2.5-flash}
qua SDK \texttt{google-genai} -- với \textit{structured outputs} ràng buộc bằng
JSON Schema để loại trừ hoàn toàn lỗi phân tích văn bản tự do từ LLM.

\subsection{Luồng xử lý}

Luồng xử lý ở mức tài liệu được tóm tắt trong Hình~\ref{fig:kpi-pipeline}.
Tệp đầu vào là JSONL đã gán nhãn ở Mục~2; tệp đầu ra là một tệp JSON cho mỗi
trang, mỗi tệp chứa một danh sách các đối tượng KPI.

\begin{figure}[H]
  \centering
  \begin{tikzpicture}[
      node distance=7mm,
      box/.style={draw, rounded corners, align=center, inner sep=4pt,
                  text width=11.5cm, minimum height=8mm},
      gpu/.style={draw, rounded corners, align=center, inner sep=4pt,
                  text width=11.5cm, minimum height=8mm, very thick},
      arrow/.style={-{Stealth[length=2.5mm]}, thick}
    ]
    \node[box] (in)   {JSONL câu đã gán nhãn ESG \\ \footnotesize (đầu ra của Mục 2)};
    \node[box, below=of in] (group) {Nhóm câu theo \texttt{(source\_pdf, page)} và sắp xếp theo \texttt{sentence\_index}};
    \node[box, below=of group] (gate) {Cổng lọc ESG: bỏ qua trang không có câu \texttt{esg=true} \\ \footnotesize (có thể tắt bằng \texttt{-{}-all-pages})};
    \node[gpu, below=of gate] (llm) {Gọi LLM Gemini 2.5 Flash với \textit{response\_schema} \\ \footnotesize (toàn bộ văn bản trang được đưa vào, không chỉ câu ESG)};
    \node[box, below=of llm] (norm) {Chuẩn hoá kết quả: \%, ép kiểu số, ép năm};
    \node[box, below=of norm] (out)  {Ghi \texttt{page\_NNN\_kpis.json} (rỗng \texttt{[]} nếu không có KPI)};
    \draw[arrow] (in) -- (group);
    \draw[arrow] (group) -- (gate);
    \draw[arrow] (gate) -- (llm);
    \draw[arrow] (llm) -- (norm);
    \draw[arrow] (norm) -- (out);
  \end{tikzpicture}
  \caption{Luồng trích xuất KPI ở mức tài liệu. Khối in đậm là lời gọi LLM có
  ràng buộc lược đồ; các bước còn lại chạy cục bộ, không tốn token.}
  \label{fig:kpi-pipeline}
\end{figure}

\subsubsection{Tái dựng văn bản trang}

JSONL đầu vào ở mức câu. Để cung cấp đủ ngữ cảnh cho LLM, mô-đun gom các dòng
thành cấu trúc \texttt{\{source\_pdf: \{page: [(sentence\_index, text,
esg), ...]\}\}}, dùng \texttt{collections.OrderedDict} để giữ thứ tự xuất hiện
của tài liệu. Văn bản mỗi trang sau đó được ghép lại theo \texttt{sentence\_index}
tăng dần (\texttt{build\_page\_text}). Quan trọng là toàn bộ trang được tái
dựng, kể cả các câu mà mô hình phân loại đã gán \textit{Neutral} -- vì con số
hỗ trợ cho một tuyên bố ESG (ví dụ chú thích bảng, mệnh đề tiếp theo) thường
nằm ở chính những câu Neutral đó; cắt bỏ chúng sẽ làm mất giá trị.

\subsubsection{Cổng lọc ESG: tiết kiệm chi phí, không hi sinh độ phủ}

Phần lớn các trang trong một báo cáo thường niên là báo cáo tài chính, ảnh, danh
sách cổ đông -- không có nội dung ESG. Ở chế độ mặc định, bước này chỉ gọi LLM
khi trang có \textbf{ít nhất một câu} với \texttt{esg=true}; các trang còn lại
được ghi sẵn dưới dạng \texttt{[]} và không tốn token. Khi cổng kích hoạt, văn
bản gửi tới LLM là \textbf{toàn bộ trang} chứ không phải tập con các câu ESG --
nghĩa là cổng chỉ quyết định \textit{có gọi hay không}, không quyết định
\textit{mô hình nhìn thấy gì}. Cờ \texttt{--all-pages} cho phép bỏ cổng, dùng
trong các thử nghiệm thu hồi (recall).

\subsubsection{Câu lệnh prompt}

Prompt được tách thành hai phần:

\begin{itemize}
  \item \textbf{System instruction} -- gán định danh (\texttt{ESG-KPI-EXTRACTOR-V2}),
        bắt buộc đầu ra JSON, \textit{ấn định cứng} các trường siêu dữ liệu
        \texttt{company}, \texttt{sector}, \texttt{page}, \texttt{doc\_name} cho
        mọi bản ghi (mô hình không có cơ hội bịa các trường này), và liệt kê
        bốn quy tắc phân loại \texttt{kind} kèm từ khoá kích hoạt song ngữ
        (\textit{baseline}: ``kể từ / năm gốc / since''; \textit{target}: ``mục
        tiêu / cam kết / hướng tới''; \textit{achieved}: ``đạt được / đã giảm /
        đã thực hiện''; \textit{projection}: ước lượng tương lai chưa cam kết).
  \item \textbf{User message} -- bơm \textit{toàn bộ bộ từ vựng KPI}
        (\texttt{\{id\}: \{definition\}} mỗi dòng, từ
        \texttt{kpi\_definitions\_construction.json}) cùng văn bản trang đã tái
        dựng đặt trong dấu \texttt{"""..."""} với số trang và tên tài liệu.
\end{itemize}

Hai quy ước có chủ ý: nếu một chỉ số không khớp định nghĩa nào, mô hình phải đặt
\texttt{kpi\_type = "other"} với \texttt{title} mô tả -- thà giữ lại quan sát có
kiểu \texttt{other} còn hơn mất thông tin; và \texttt{source\_id} được mẫu hoá
\texttt{"\{doc\}\_\{page\}\_\{index\}"} để mỗi quan sát có một khoá ổn định, sau
này dùng làm \texttt{identity\_keys} khi tải vào đồ thị.

\subsubsection{Sinh có ràng buộc lược đồ}

Lời gọi LLM truyền vào \texttt{response\_mime\_type = "application/json"} cùng
một \texttt{response\_schema} theo \textit{phương ngữ OpenAPI-3 của Gemini} (các
trường có thể \texttt{null} được khai báo bằng \texttt{"nullable": true} chứ
không phải union, và không dùng \texttt{additionalProperties}). Các trường có
thể trống về mặt nội dung (\texttt{value}, \texttt{unit}, \texttt{year},
\texttt{target\_year}, \texttt{baseline\_year}) vẫn được liệt kê trong
\texttt{required} -- nghĩa là mô hình \textbf{bắt buộc phát ra} khoá ấy, dù giá
trị là \texttt{null}. Điều này loại bỏ hiện tượng thiếu khoá ngầm và giúp bước
hậu xử lý phân biệt được ``không biết'' và ``mô hình quên''. Tham số
\texttt{temperature = 0} bảo đảm tính tái lập;
\texttt{max\_output\_tokens = 8000} đủ rộng cho các trang nhiều KPI.

\subsubsection{Phòng vệ trước các trạng thái dừng bất thường}

Ngay cả khi đã ràng buộc lược đồ, Gemini vẫn có thể kết thúc vì lý do khác
\texttt{STOP}: bộ lọc an toàn, lặp văn bản, ... Mã chỉ chấp nhận
\texttt{finish\_reason $\in$ \{STOP, MAX\_TOKENS\}}; mọi trạng thái khác sẽ ghi
cảnh báo và trả về \texttt{[]} cho trang đó, không ném ngoại lệ -- nhờ vậy một
trang bị chặn bởi bộ lọc nội dung không làm sụp đổ cả tài liệu.

\subsubsection{Chuẩn hoá hậu xử lý}

Mô hình thi thoảng trả về chuỗi \texttt{"15\%"} (vì văn bản gốc viết đúng như
vậy). Hàm \texttt{normalize\_kpi\_response} bóc dấu \texttt{\%} cuối chuỗi,
nâng lên \texttt{unit} nếu \texttt{unit} đang trống, rồi ép \texttt{value} sang
\texttt{float}; đồng thời ép các trường năm dạng chuỗi 4 chữ số (\texttt{"2025"})
về \texttt{int}. Đây là biến đổi duy nhất sau khi LLM trả lời -- không bịa,
không suy diễn giá trị.

\subsubsection{Song song hoá và khả năng tiếp tục (resumable)}

Các trang của cùng một tài liệu được xử lý song song bằng
\texttt{ThreadPoolExecutor} với \texttt{-{}-max-workers} luồng (mặc định 4); chọn
thread (không phải process) là phù hợp vì nghẽn cổ chai nằm ở I/O mạng tới
API Gemini, và SDK của Google là thread-safe. Trước mỗi lần gửi, worker kiểm tra
tệp đầu ra \texttt{page\_NNN\_kpis.json} đã tồn tại hay chưa; nếu có thì đọc
lại và bỏ qua. Cơ chế này khiến chạy lại an toàn: một lượt chạy 1.000 trang bị
ngắt ở trang 700, khi chạy lại chỉ tốn token cho 300 trang còn lại. Trang xảy
ra lỗi được ghi nhận và \textbf{không} ghi tệp đầu ra (để lượt chạy sau thử
lại), thay vì cache lỗi thành \texttt{[]}.

\subsection{Cấu trúc bản ghi kết quả}

Mỗi tệp \texttt{page\_NNN\_kpis.json} là một danh sách các đối tượng KPI có dạng
như Bảng~\ref{tab:kpi-record}. Một KPI có thể chứa nhiều \texttt{observations}
khi cùng một chỉ số được nêu ở nhiều thời điểm (kỳ gốc + mục tiêu + đã đạt).

\begin{table}[H]
  \centering
  \caption{Các trường chính trong bản ghi KPI đầu ra.}
  \label{tab:kpi-record}
  \begin{tabularx}{\linewidth}{l X}
    \hline
    \textbf{Trường} & \textbf{Ý nghĩa} \\
    \hline
    \texttt{kpi\_type}   & Mã KPI trong bộ từ vựng kiểm soát; \texttt{"other"} nếu không khớp. \\
    \texttt{title}       & Nhãn ngắn mô tả chỉ số. \\
    \texttt{observations}& Danh sách quan sát; mỗi quan sát có \texttt{value, unit, kind, direction, year, target\_year, baseline\_year, source\_id, snippet}. \\
    \texttt{kind}        & Một trong \{\texttt{baseline}, \texttt{target}, \texttt{achieved}, \texttt{projection}\}. \\
    \texttt{direction}   & Một trong \{\texttt{absolute}, \texttt{reduction}, \texttt{increase}\}. \\
    \texttt{source\_id}  & Khoá theo dạng \texttt{"\{doc\}\_\{page\}\_\{index\}"} để truy vết và chống trùng. \\
    \texttt{snippet}     & Đoạn trích nguồn $\le$ 160 ký tự để kiểm chứng. \\
    \texttt{page, doc\_name, company, sector} & Siêu dữ liệu cố định cho mọi bản ghi của trang. \\
    \hline
  \end{tabularx}
\end{table}

Tổ chức một tệp cho mỗi trang phục vụ ba mục đích: (i) gắn số trang vào mọi bản
ghi để có thể truy ngược chính xác về tài liệu nguồn; (ii) làm cho quy trình
\textit{tiếp tục được} khi bị gián đoạn; và (iii) giới hạn ngữ cảnh mỗi lời gọi
LLM ở mức một trang, giúp độ trễ và lượng token có thể dự đoán.

\subsection{So sánh với pipeline tham chiếu}

Mô-đun này khởi đầu từ việc cổng (port) lại bước
\texttt{1-kpi-extraction.py} của dự án EmeraldMind, giữ nguyên định dạng bản
ghi, hàm chuẩn hoá và khung prompt. Những khác biệt có chủ đích được nêu trong
Bảng~\ref{tab:kpi-vs-emerald}.

\begin{table}[H]
  \centering
  \caption{Khác biệt giữa mô-đun trích xuất KPI của dự án và bước 1 của EmeraldMind.}
  \label{tab:kpi-vs-emerald}
  \begin{tabularx}{\linewidth}{l X X}
    \hline
    \textbf{Khía cạnh} & \textbf{EmeraldMind bước 1} & \textbf{Dự án này} \\
    \hline
    Đầu vào            & Tệp PDF gốc & JSONL câu đã trích xuất \& gán nhãn \\
    Bộ từ vựng KPI     & Đa ngành, có FAISS để chọn ngành & Một ngành (Xây dựng -- VLXD -- BĐS), không cần FAISS \\
    Phát hiện ngành    & Heuristic theo tài liệu & Hằng số cứng \\
    LLM                & Gemini với 7 khoá luân phiên (để né rate limit) & Gemini 2.5 Flash, một khoá \\
    Định dạng đầu ra   & JSON lẫn trong văn bản, phải bóc khung mã & Structured outputs (\texttt{response\_schema}) \\
    Cổng xử lý         & Mọi trang & Mặc định chỉ trang có câu ESG, \texttt{-{}-all-pages} để bỏ cổng \\
    \hline
  \end{tabularx}
\end{table}

\medskip
\noindent\textbf{Đầu ra cuối cùng của giai đoạn này} là kho bản ghi KPI có cấu
trúc, neo về bộ từ vựng kiểm soát và truy vết về tài liệu/trang nguồn. Đây
chính là tập dữ liệu sẽ được nạp vào nút \texttt{KPIObservation} của đồ thị tri
thức trong bước xây dựng GraphRAG tiếp theo.

\end{document}
